% ch-a3-03 —— 第 3 篇 Fargues–Scholze，PDF 162–170，印刷页 149–157
\ysrc{149}
其中用到公式
\[
\mathcal{H}ck_{G,\mathrm{Div}^{d}}=[L^{+}G\backslash LG/L^{+}G],
\]
这里 $LG$ 与 $L^{+}G$ 分别是 $\mathrm{Div}^{d}$ 上的圈群与正圈群。使扭叠积（twisted stacky product）得以定义的那个 $L^{+}G$ 在 $L^{+}G\backslash LG\times LG/L^{+}G$ 上的作用，对 $x,y\in LG$ 与 $g\in L^{+}G$ 由 $g.([x],[y])=([xg^{-1}],[gy])$ 给出。态射 $a$ 由 $a([([x],[y])])=([x],[y])$ 给出，态射 $b$ 由 $b([([x],[y])])=[xy]$ 给出。

沿 $(\mathrm{Div}^{1})^{I}\to\mathrm{Div}^{d}$ 拉回到 $(\mathrm{Div}^{1})^{I}$ 之后，上述构造给出两个态射 $a',b'$；于是对 $A_{1},A_{2}\in D_{\mathrm{ét}}(\mathcal{H}ck^{I}_{G},\Lambda)^{bd}$（上标「bounded」意为具有拟紧支撑），可以定义二者的卷积
\[
A_{1}\star A_{2}=Rb'_{*}a'^{*}(A_{1}\boxtimes A_{2})\in D_{\mathrm{ét}}(\mathcal{H}ck^{I}_{G},\Lambda)^{bd}.
\]
基本结果如下。

\begin{thm}{4.4}
运算 $\star$ 保持 $\mathrm{Sat}^{I}_{G}(\Lambda)$，并定义卷积积
\[
\star:\mathrm{Sat}^{I}_{G}(\Lambda)\times\mathrm{Sat}^{I}_{G}(\Lambda)\longrightarrow\mathrm{Sat}^{I}_{G}(\Lambda).
\]
\label{thm:a3-4-4}
\end{thm}

\subsubsection{融合（[9, Chapter VI.9]）}
卷积积的主要问题在于证明它在 Satake 范畴上定义对称幺半结构。这一点借助融合积（fusion product）、仿照 [20] 实现。这里本质地用到 Satake 范畴定义中的 ULA 性质（universally locally acyclic，普遍局部无环）。更确切地说，设 $I=I_{1}\sqcup\cdots\sqcup I_{k}$。按如下方式定义一个幺半函子
\[
\mathrm{Sat}^{I_{1}}_{G}(\Lambda)\times\cdots\times\mathrm{Sat}^{I_{k}}_{G}(\Lambda)\longrightarrow\mathrm{Sat}^{I}_{G}(\Lambda).
\]
为此考虑开子集
\[
j:(\mathrm{Div}^{1})^{I;I_{1},\dots,I_{k}}\hookrightarrow(\mathrm{Div}^{1})^{I},
\]
其中当 $i,i'\in I$ 落在不同的 $I_{j}$ 中时 $x_{i}\ne x_{i'}$（作为曲线上的点）。与定义 $\mathrm{Sat}^{I}_{G}$ 的方式相同，定义
\[
\mathrm{Sat}^{I;I_{1},\dots,I_{k}}_{G}(\Lambda)\subset D_{\mathrm{ét}}\bigl(\mathcal{H}ck^{I}_{G}\times_{(\mathrm{Div}^{1})^{I}}(\mathrm{Div}^{1})^{I;I_{1},\dots,I_{k}},\Lambda\bigr)
\]
为有界 ULA 平坦反常层构成的全子范畴。

\begin{thm}{4.5}
开浸入 $j:\mathcal{H}ck^{I}_{G}\times_{(\mathrm{Div}^{1})^{I}}(\mathrm{Div}^{1})^{I;I_{1},\dots,I_{k}}\hookrightarrow\mathcal{H}ck^{I}_{G}$ 诱导全忠实函子
\[
j^{*}:\mathrm{Sat}^{I}_{G}(\Lambda)\longrightarrow\mathrm{Sat}^{I;I_{1},\dots,I_{k}}_{G}(\Lambda).
\]
\label{thm:a3-4-5}
\end{thm}

此处全忠实性至关重要，它在后文强制交换约束成立。下一个结果如下。对 $1\le r\le k$ 有态射
\[
\mathcal{H}ck^{I}_{G}\times_{(\mathrm{Div}^{1})^{I}}(\mathrm{Div}^{1})^{I;I_{1},\dots,I_{k}}\longrightarrow\mathcal{H}ck^{I_{r}}_{G}.
\]
由此可以对 $A_{1}\in\mathrm{Sat}^{I_{1}}_{G}(\Lambda),\cdots,A_{k}\in\mathrm{Sat}^{I_{k}}_{G}(\Lambda)$ 定义
\[
A_{1}\boxtimes\cdots\boxtimes A_{k}\in\mathrm{Sat}^{I;I_{1},\dots,I_{k}}_{G}(\Lambda).
\]

\begin{thm}{4.6}
外张量积
\[
\mathrm{Sat}^{I_{1}}_{G}(\Lambda)\times\cdots\times\mathrm{Sat}^{I_{k}}_{G}(\Lambda)\longrightarrow\mathrm{Sat}^{I;I_{1},\dots,I_{k}}_{G}(\Lambda)
\]
的像，经定理 4.5 的全忠实函子 $j^{*}$ 落在 $\mathrm{Sat}^{I}_{G}(\Lambda)$ 中。
\label{thm:a3-4-6}
\end{thm}

\ysrc{150}
定理 4.6 的证明依赖一个临时的辅助构造：局部 Hecke 叠（local Hecke stack）
\[
\mathcal{H}ck^{I,I_{1},\dots,I_{k}}\longrightarrow(\mathrm{Div}^{1})^{I},
\]
它带有投影
\[
p_{j}:\mathcal{H}ck^{I,I_{1},\dots,I_{k}}\longrightarrow\mathcal{H}ck^{I_{j}},\qquad j=1,\dots,k,
\]
以及一个态射
\[
m:\mathcal{H}ck^{I,I_{1},\dots,I_{k}}\longrightarrow\mathcal{H}ck^{I}.
\]
其定义如下：$\mathcal{H}ck^{I,I_{1},\dots,I_{k}}(S)$ 指的是下述数据：
\begin{itemize}
\item 一族次数为 $1$ 的相对 Cartier 除子 $(D_{i})_{i\in I}$；
\item $X_{S}$ 上的一族 $G$-丛 $\mathcal{E}_{0},\dots,\mathcal{E}_{k}$；
\item 同构
\[
\mathcal{E}_{j-1}|_{X_{S}\setminus\coprod_{i\in I_{j}}D_{i}}\xrightarrow{\ \sim\ }\mathcal{E}_{j}|_{X_{S}\setminus\coprod_{i\in I_{j}}D_{i}},\qquad 1\le j\le k,
\]
\end{itemize}
而这些同构沿 Cartier 除子亚纯。投影 $p_{j}$ 把上述数据送到二元组 $(\mathcal{E}_{j-1},\mathcal{E}_{j})$，连同上面第 $j$ 个同构给出的修正。态射 $m$ 把上述数据送到二元组 $(\mathcal{E}_{0},\mathcal{E}_{k})$，连同把 $k$ 个修正限制到 $X_{S}\setminus\coprod_{i\in I}D_{i}$ 后复合所得的修正。于是可以验证
\[
B=Rm_{*}\bigl(p_{1}^{*}A_{1}\otimes^{L}_{\Lambda}\cdots\otimes^{L}_{\Lambda}p_{k}^{*}A_{k}\bigr)\in D_{\mathrm{ét}}(\mathcal{H}ck^{I},\Lambda)
\]
实际上落在 $\mathrm{Sat}^{I}_{G}(\Lambda)$ 中。由于它沿 $j$ 的拉回是 $A_{1}\boxtimes\cdots\boxtimes A_{k}$，定理 4.6 得证。

现在可以构造融合积。对应 $I\mapsto\mathrm{Sat}^{I}_{G}(\Lambda)$ 对 $I$ 具有函子性：对有限集的映射 $I\to J$，有函子 $\mathrm{Sat}^{I}_{G}(\Lambda)\to\mathrm{Sat}^{J}_{G}(\Lambda)$。于是可以借助映射 $I\sqcup\cdots\sqcup I\to I$ 作复合
\[
\mathrm{Sat}^{I}_{G}(\Lambda)\times\cdots\times\mathrm{Sat}^{I}_{G}(\Lambda)\longrightarrow\mathrm{Sat}^{I\sqcup\cdots\sqcup I}_{G}(\Lambda)\longrightarrow\mathrm{Sat}^{I}_{G}(\Lambda).
\]
这就把融合积 $\ast$ 定义成卷积积 $\star$ 的对称幺半细化（它确实细化了卷积积，这一点是上述定理证明的一部分）。由此得到函子
\[
I\longmapsto(\mathrm{Sat}^{I}_{G}(\Lambda),\ast)
\]
从有限集到对称幺半范畴。此外还可以证明，函子
\[
F^{I}:\mathrm{Sat}^{I}_{G}(\Lambda)\longrightarrow\Rep_{\Lambda}(W^{I}_{E})
\]
是对称幺半的。

\subsection{Tannakian 重构（[9, Chapter VI.10]）}
设 ${}^{L}G=\widehat{G}\rtimes W_{E}$ 是 $G$ 在 $\Lambda$ 上的 Langlands 对偶。可以把 $\Rep_{{}^{L}G}(\Lambda)$ 定义为 ${}^{L}G$ 在有限型投射 $\Lambda$-模上的表示构成的范畴。这里要求表示限制到 $\widehat{G}$ 时是代数的，且在 $W_{E}$ 的某个开子群上平凡。对应
\[
I\longmapsto\Rep_{\Lambda}({}^{L}G^{I})
\]
是有限集到对称幺半范畴的函子。

\ysrc{151}
\begin{thm}{4.7}
对 $\mathbb{Z}_{\ell}[\sqrt{q}]$-代数 $\Lambda$，存在从有限集到对称幺半范畴的函子之间的等价
\[
\Rep_{\Lambda}({}^{L}G^{I})\xrightarrow{\ \sim\ }\mathrm{Sat}^{I}_{G}(\Lambda).
\]
通过该等价，函子 $F^{I}:\mathrm{Sat}^{I}_{G}(\Lambda)\to\Rep_{\Lambda}W^{I}_{E}$ 等同于限制到 $W^{I}_{E}$。由拉回函子 $\mathcal{H}ck^{I}_{G}\to(\mathrm{Div}^{1})^{I}$ 诱导的自然包含 $\Rep_{\Lambda}({}^{L}G^{I})\to\mathrm{Sat}^{I}_{G}$，对应于沿 $({}^{L}G^{I})\to W^{I}_{E}$ 的拉回。
\label{thm:a3-4-7}
\end{thm}

\subsection{我们使用的工具}
\subsubsection{双曲局部化}
我们用到的首要工具是双曲局部化（hyperbolic localization，[9, Chapter IV.6]），它沿用 [21] 的工作。定理 4.1 的证明依赖于此。双曲局部化使我们能证明下述结论，见 [9, Chapter VI.3]。设 $G$ 分裂。取 $T\subset B$ 为 $G$ 中某个 Borel 子群内的极大环面。考虑与 $B$ 相伴的仿射 Grassmann 簇
\[
\mathrm{Gr}_{B,\mathrm{Div}^{d}}\longrightarrow\mathrm{Div}^{d}.
\]
商态射 $B\to T$ 诱导态射
\[
\mathrm{Gr}_{B,\mathrm{Div}^{d}}\longrightarrow\mathrm{Gr}_{T,\mathrm{Div}^{d}}=\coprod_{X_{*}(T)}\mathrm{Div}^{d}.
\]
由此通过拉回得到分解
\[
\mathrm{Gr}_{B,\mathrm{Div}^{d}}=\coprod_{\lambda\in X_{*}(T)}\mathrm{Gr}^{\lambda}_{B,\mathrm{Div}^{d}}.
\]
于是有自然态射
\[
\mathrm{Gr}_{B,\mathrm{Div}^{d}}\longrightarrow\mathrm{Gr}_{G,\mathrm{Div}^{d}},
\]
它在几何点上双射，且在每片 $\mathrm{Gr}^{\lambda}_{B,\mathrm{Div}^{d}}$（$\lambda\in X_{*}(T)$）上是局部闭浸入。像 $\mathrm{Gr}^{\lambda}_{B,\mathrm{Div}^{d}}\hookrightarrow\mathrm{Gr}_{G,\mathrm{Div}^{d}}$ 就是所谓的半无限轨道（semi-infinite orbit）。令
\[
\begin{tikzcd}[column sep=large,row sep=large]
\mathrm{Gr}_{B,\mathrm{Div}^{d}} \arrow[r, "q"] \arrow[d, "p"'] & \mathrm{Gr}_{G,\mathrm{Div}^{d}} \\
\mathrm{Gr}_{T,\mathrm{Div}^{d}} &
\end{tikzcd}
\]
于是可以定义常数项函子（constant term functor）
\[
\mathrm{CT}_{B}:D_{\mathrm{ét}}(\mathcal{H}ck^{I}_{G},\Lambda)^{bd}\longrightarrow D_{\mathrm{ét}}(\mathcal{H}ck^{I}_{T},\Lambda)^{bd}
\]
为 $Rp_{!}q^{*}$。下述结果本质地用到双曲局部化。关于这一切可参看 [9, Chapter VI.3]。移位「deg」是显式的局部常值函数 $\mathrm{Gr}_{T,\mathrm{Div}^{d}}\to\mathbb{Z}$，其存在只是为了保证反常层的常数项仍是反常层。

\begin{thm}{4.8}
常数项函子满足
\[
\mathrm{CT}_{B}[\deg]:\mathrm{Sat}^{I}_{G}(\Lambda)\longrightarrow\mathrm{Sat}^{I}_{T}(\Lambda),
\]
并且
\[
R\pi_{G*}=R\pi_{T*}\circ\mathrm{CT}_{B}[\deg]:\mathrm{Sat}^{I}_{G}(\Lambda)\longrightarrow\LocSys((\mathrm{Div}^{1})^{I},\Lambda).
\]
\label{thm:a3-4-8}
\end{thm}

\ysrc{152}
这是我们分析范畴 $\mathrm{Sat}^{I}_{G}(\Lambda)$ 的主要工具。事实上，$\mathrm{CT}_{B}[\deg]$ 把限制函子 $\Rep_{\Lambda}({}^{L}G^{I})\to\Rep_{\Lambda}({}^{L}T^{I})$ 几何化，这解释了它在 Satake 等价证明中的重要性。

\subsubsection{退化到 Witt 向量 Grassmann 簇}
设 $G$ 分裂。现在把它看作 $\mathcal{O}_{E}$ 上的约化群。对 $S\in\Perf_{\mathbb{F}_{q}}$，存在一个 sous-完美胚 空间\tn{底本作“sous-完美胚”，系法文“sous-perfectoïde”的残留写法，当指次完美胚（sub-完美胚）空间；上下文推断，无外部依据。}
\[
\mathcal{Y}_{S}\longrightarrow\Spa(\mathcal{O}_{\breve{E}})
\]
使得
\[
\mathcal{Y}_{S}^{\diamond}=S\times\mathrm{Spd}(\mathcal{O}_{\breve{E}}).
\]
有 $\{\pi\ne0\}=Y_{S}\subset\mathcal{Y}_{S}\supset S=\{\pi=0\}$。可以定义
\[
\mathrm{Div}^{1}_{\mathcal{Y}}=\mathrm{Spd}(\mathcal{O}_{\breve{E}})
\]
以及相应的 v-层
\[
\mathrm{Gr}_{G,\mathrm{Div}^{1}_{\mathcal{Y}}}\longrightarrow\mathrm{Div}^{1}_{\mathcal{Y}}.
\]
设 $\mathrm{Gr}^{\mathrm{Witt}}_{G,\mathbb{F}_{q}}\to\Spec(\mathbb{F}_{q})$ 是 Zhu 的 Witt 向量仿射 Grassmann 簇（[28], [3]）。有
\[
\mathrm{Gr}_{G,\mathrm{Div}^{1}_{\mathcal{Y}}}\times_{\mathrm{Div}^{1}_{\mathcal{Y}}}\mathrm{Spd}(\mathbb{F}_{q})=\bigl(\mathrm{Gr}^{\mathrm{Witt}}_{G,\mathbb{F}_{q}}\bigr)^{\diamond}.
\]
在重构定理 4.7 的证明中，这一点用来证明：对 $\mu\in X_{*}(T)^{+}$，若 $j_{\mu}:\mathcal{H}ck_{G,\mu}\hookrightarrow\mathcal{H}ck_{G}$ 是相应仿射 Schubert 胞腔（维数为 $d_{\mu}$）的包含，则
\[
{}^{p}j_{\mu!}\mathbb{Q}_{\ell}[d_{\mu}]\longrightarrow{}^{p}Rj_{\mu*}\mathbb{Q}_{\ell}[d_{\mu}]
\]
是同构。这一结论可以转移到 $\mathrm{Gr}^{\mathrm{Witt}}_{G,\mathbb{F}_{q}}$ 上的同类命题，那里的证明最终用到 Demazure 消解上的分解定理，见 [9, Proposition VI.7.5]。

\section{Langlands 参数与谱作用}
\subsection{Langlands 参数的模叠（[9, Chapter VIII.1], [5]）}
\subsubsection{存在性与奇点}
可以把任何 $\mathbb{Z}_{\ell}$-代数 $\Lambda$ 看作凝聚环（condensed ring），方式是令 $\Lambda:=\Lambda^{\mathrm{disc}}\otimes_{\mathbb{Z}_{\ell}^{\mathrm{disc}}}\mathbb{Z}_{\ell}$。定义 $Z^{1}(W_{E},\widehat{G})$ 为 $\mathbb{Z}_{\ell}$-代数上的函子，它在 $\Lambda$ 上的取值是凝聚 1-上闭链 $W_{E}\to\widehat{G}(\Lambda)$，其中 $\widehat{G}(\Lambda)$ 自然地是凝聚群。$\widehat{G}$ 通过共轭作用在 $Z^{1}(W_{E},\widehat{G})$ 上。关于它的第一个主要结果如下。

\begin{thm}{5.1}
函子 $Z^{1}(W_{E},\widehat{G})$ 可由一个 $\mathbb{Z}_{\ell}$-概形表示，后者是有限型仿射概形的无限无交并。此外，代数叠
\[
[Z^{1}(W_{E},\widehat{G})/\widehat{G}]
\]
是 $\Spec(\mathbb{Z}_{\ell})$ 上维数为 0 的局部完全交。
\label{thm:a3-5-1}
\end{thm}

我们采用记号
\[
\LocSys_{\widehat{G}}=[Z^{1}(W_{E},\widehat{G})/\widehat{G}].
\]

\begin{rem}{5.2}
我们的参数模叠是维数为 0 的局部完全交。我们不需要像紧 Riemann 面情形那样把它提升为导出叠——在那种情形，朴素的非导出局部系统模空间不具有上述性质。我们可以考虑自己模空间的导出版本，但最终会证明它实际上是非导出的。因此从这一观点看，我们的处境比紧 Riemann 面上的几何 Langlands 纲领更好。
\label{rem:a3-5-2}
\end{rem}

\ysrc{153}
这一切都源于 Frobenius 的在场及其对温和惯性（moderate inertia）的作用，也源于 $\widehat{\mathbb{Z}}{}^{p}(1)$ 的在场——它作为 $I_{E}/P_{E}$ 出现。

\subsubsection{粗模空间}
粗模空间是
\[
Z^{1}(W_{E},\widehat{G})\sslash\widehat{G},
\]
它是有限型 $\mathbb{Z}_{\ell}$-仿射概形的无限无交并。其函数环为
\[
\mathcal{O}(Z^{1}(W_{E},\widehat{G}))^{\widehat{G}}.
\]
其几何点可以借助几何不变量理论、用 Hilbert–Mumford–Kempf 数值判别法描述，见 [9, Chapter VIII.3.1]。

\begin{thm}{5.3}
设 $L$ 是 $\mathbb{Z}_{\ell}$ 上的代数闭域。则 $\LocSys_{\widehat{G}}$ 的粗模空间的 $L$-点，由半单 Langlands 参数 $W_{E}\to\widehat{G}(L)\rtimes W_{E}$ 的 $\widehat{G}(L)$-共轭类给出。
\label{thm:a3-5-3}
\end{thm}

按定义，$\widehat{G}_{L}\rtimes W_{E}$ 的抛物子群指的是满射到 $W_{E}$、且与 $\widehat{G}$ 的交为 $\widehat{G}$ 的抛物子群的子群。在 $\widehat{G}(L)$-共轭下，这些恰好是 $\widehat{P}_{L}\rtimes W_{E}$，其中 $P$ 是 $G^{*}$——$G$ 的拟分裂内型——的抛物子群。Levi 子群的情形同理。回忆半单 Langlands 参数指的是参数 $\varphi:W_{E}\to\widehat{G}(L)\rtimes W_{E}$，使得若其像含于某个抛物子群，则其像含于该抛物子群的某个 Levi 子群。

若 $L=\mathbb{Q}_{\ell}$ 或 $\overline{\mathbb{Q}}_{\ell}$，参数由 Weil–Deligne 群的表示给出。这样的表示由一对 $(\rho,N)$ 组成，其中 $\rho:W_{E}\to\widehat{G}(L)$ 是在 $I_{E}$ 的某个开子群上平凡的参数，而 $N\in\operatorname{Lie}\widehat{G}\otimes L$ 满足
\[
\mathrm{Ad}(\rho(\tau)).N=q^{v(\tau)}N
\]
对一切 $\tau\in W_{E}$ 成立，其中 $\tau$ 在剩余域上通过 $\mathrm{Frob}_{q}^{v(\tau)}$ 作用。于是若 $\varphi$ 对应于 $(\rho,N)$，则 $\varphi$ 半单当且仅当 $N=0$ 且 $\varphi$ 是 Frobenius 半单的。

\subsubsection{无穷小性质（[9, Chapter VII.2]）}
最后看这个模空间的无穷小性质，也就是余切复形
\[
\mathbb{L}_{\LocSys_{\widehat{G}}/\mathbb{Z}_{\ell}}\in\Perf^{[-1,1]}(\LocSys_{\widehat{G}}),
\]
这里用到它是 $\mathbb{Z}_{\ell}$ 上的局部完全交。有如下描述。

\begin{thm}{5.4}
设 $x:\Spec(\Lambda)\to\LocSys_{\widehat{G}}$ 对应于 $\varphi:W_{E}\to\widehat{G}(\Lambda)$。则
\[
x^{*}\mathbb{L}^{\vee}_{\LocSys_{\widehat{G}}/\mathbb{Z}_{\ell}}=R\Gamma\bigl(W_{E},(\widehat{\mathfrak{g}}^{*}\otimes_{\mathbb{Z}_{\ell}}\Lambda)_{\varphi}(1)\bigr)[1],
\]
其中 $(\widehat{\mathfrak{g}}^{*}\otimes_{\mathbb{Z}_{\ell}}\Lambda)_{\varphi}$ 指 $\widehat{\mathfrak{g}}^{*}\otimes_{\mathbb{Z}_{\ell}}\Lambda$ 带上由 $\mathrm{Ad}\circ\varphi$ 导出的 $W_{E}$ 的扭作用。
\label{thm:a3-5-4}
\end{thm}

后文讨论 Arinkin–Gaitsgory 奇异支撑条件时，需要这一显式描述。令
\[
\Sing_{\LocSys_{\widehat{G}}/\mathbb{Z}_{\ell}}:=\mathcal{H}^{-1}\bigl(\mathbb{L}_{\LocSys_{\widehat{G}}/\mathbb{Z}_{\ell}}\bigr)\longrightarrow\LocSys_{\widehat{G}}
\]

\ysrc{154}
称为奇点叠。这里 $\mathcal{H}^{-1}$ 指：这是 $\LocSys_{\widehat{G}}$ 上的群概形，它表示把带有态射 $f:S\to\LocSys_{\widehat{G}}$ 的概形 $S$ 送到 $H^{-1}(S,f^{*}\mathbb{L}_{\LocSys_{\widehat{G}}/\mathbb{Z}_{\ell}})$ 的函子。在这类 $S$ 上局部地，若把 $f^{*}\mathbb{L}_{\LocSys_{\widehat{G}}/\mathbb{Z}_{\ell}}\simeq[\mathcal{E}^{-1}\to\mathcal{E}^{0}\to\mathcal{E}^{1}]$ 写成向量丛 $\mathcal{E}^{-1},\mathcal{E}^{0},\mathcal{E}^{1}$ 构成的复形，则它由 $\ker(\mathbb{V}(\mathcal{E}^{-1})\to\mathbb{V}(\mathcal{E}^{0}))$ 表示。该群概形带有 $\mathbb{G}_{m}$ 的作用，因而可以谈论其中的齐性 Zariski 闭子集。

\begin{cor}{5.5}
对 $x:\Spec(\Lambda)\to\LocSys_{\widehat{G}}$，有
\[
x^{*}\Sing_{\LocSys_{\widehat{G}}/\mathbb{Z}_{\ell}}=\bigl\{v\in\widehat{\mathfrak{g}}^{*}\otimes_{\mathbb{Z}_{\ell}}\Lambda\mid\forall\tau\in W_{E},\ q^{v(\tau)}\mathrm{Ad}(\varphi(\tau))(\tau.v)=v\bigr\}.
\]
\label{cor:a3-5-5}
\end{cor}

因此有自然嵌入
\[
\Sing_{\LocSys_{\widehat{G}}/\mathbb{Z}_{\ell}}\hookrightarrow[\widehat{\mathfrak{g}}^{*}/\widehat{G}]\times_{B\widehat{G}}\LocSys_{\widehat{G}/\mathbb{Z}_{\ell}}.
\]
设 $\mathcal{N}^{*}_{\widehat{G}}\subset\mathfrak{g}^{*}$ 是幂零锥。由推论 5.5 推出，将 $\ell$ 可逆化之后
\[
\Sing_{\LocSys_{\widehat{G}}/\mathbb{Z}_{\ell}}\otimes_{\mathbb{Z}_{\ell}}\mathbb{Q}_{\ell}\hookrightarrow\bigl[\mathcal{N}^{*}_{\widehat{G}_{\mathbb{Q}_{\ell}}}/\widehat{G}_{\mathbb{Q}_{\ell}}\bigr]\times_{B\widehat{G}_{\mathbb{Q}_{\ell}}}\LocSys_{\widehat{G}_{\mathbb{Q}_{\ell}}/\mathbb{Q}_{\ell}}.
\]

\begin{rem}{5.6}
（承接注 5.2）作为推论，将 $\ell$ 可逆化之后，Arinkin–Gaitsgory 奇异支撑条件（见 §6.1）自动成立。这同样是紧 Riemann 面上几何 Langlands 纲领情形下的简化。原因依然是 Frobenius 的在场：正如关系式 $\varphi N\varphi^{-1}=qN$ 蕴含单值算子 $N$ 自动幂零。
\label{rem:a3-5-6}
\end{rem}

\subsection{Hecke 作用（[9, Chapter IX.2]）}
\subsubsection{定义}
这里假设 $\Lambda$ 是扭 $\mathbb{Z}_{\ell}[\sqrt{q}]$-代数，以简化陈述。对有限集 $I$ 有图表
\[
\begin{tikzcd}[column sep=small,row sep=large]
& & \mathrm{Hck}^{I}_{G} \arrow[rr] \arrow[dll, "p_{1}"'] \arrow[dr, "p_{2}"] & & \mathcal{H}ck^{I}_{G} \arrow[dr] & \\
\mathrm{Bun}_{G} & & & \mathrm{Bun}_{G}\times(\mathrm{Div}^{1})^{I} \arrow[rr] & & (\mathrm{Div}^{1})^{I}
\end{tikzcd}
\]
其中 $\mathrm{Hck}^{I}_{G}$ 是整体 Hecke 叠，它被送到 §4.1 的局部 Hecke 叠 $\mathcal{H}ck^{I}_{G}$。利用几何 Satake 对应，沿 $\mathrm{Hck}^{I}_{G}\to\mathcal{H}ck^{I}_{G}$ 拉回，对每个 $V\in\Rep_{{}^{L}G^{I}}(\Lambda)$ 得到一个复形
\[
\mathcal{S}_{V}\in D_{\mathrm{lis}}(\mathrm{Hck}^{I}_{G},\Lambda).
\]
这定义了 Hecke 作用
\[
T_{V}=Rp_{2!}\bigl(p_{1}^{*}(-)\otimes^{L}_{\Lambda}\mathcal{S}_{V}\bigr):D_{\mathrm{lis}}(\Bun_{G},\Lambda)\longrightarrow D_{\mathrm{lis}}(\Bun_{G}\times(\mathrm{Div}^{1})^{I},\Lambda).
\]
利用 Drinfeld 引理（定理 4.2）得到函子
\[
T_{V}:D_{\mathrm{lis}}(\Bun_{G},\Lambda)\longrightarrow D_{\mathrm{lis}}(\Bun_{G}\times[*/(W^{I}_{E})],\Lambda).
\]
当 $\Lambda$ 不是扭模时，多费些功夫上述构造仍然成立。

\ysrc{155}
\subsubsection{因子化性质}
(1) Satake 对应的性质表明，$T_{V}$ 的构造对 $I$ 具有函子性：若 $\alpha:I\to J$ 且 $V\in\Rep_{{}^{L}G^{I}}(\Lambda)$，则它定义了 $\alpha_{*}V\in\Rep_{{}^{L}G^{J}}(\Lambda)$，而 $T_{\alpha_{*}V}$ 是 $T_{V}$ 与函子
\[
D_{\mathrm{lis}}(\Bun_{G}\times[*/(W^{I}_{E})],\Lambda)\longrightarrow D_{\mathrm{lis}}(\Bun_{G}\times[*/(W^{J}_{E})],\Lambda)
\]
的复合，后者是由 $\alpha^{*}:(\mathrm{Div}^{1})^{J}\to(\mathrm{Div}^{1})^{I}$ 导出的拉回，
\[
\begin{tikzcd}[column sep=large,row sep=large]
D_{\mathrm{lis}}(\Bun_{G},\Lambda) \arrow[dr, "T_{\alpha^{*}V}"'] \arrow[d, "T_{V}"'] & \\
D_{\mathrm{lis}}(\Bun_{G}\times[*/(W^{I}_{E})],\Lambda) \arrow[r, "(\mathrm{Id}\times\alpha^{*})^{*}"'] & D_{\mathrm{lis}}(\Bun_{G}\times[*/(W^{J}_{E})],\Lambda).
\end{tikzcd}
\]

(2) 若 $\rho\in\Rep_{W^{I}_{E}}(\Lambda)$，它定义 $[*/(W^{I}_{E})]$ 上的局部系统 $\mathcal{F}_{\rho}$。于是有等同
\[
T_{V\otimes\rho}=\bigl[(-)\otimes^{L}_{\Lambda}(\Lambda\boxtimes\mathcal{F}_{\rho})\bigr]\circ T_{V}.
\]

(3) $T_{V_{1}\otimes V_{2}}$ 由 $T_{V_{1}}(-)\otimes^{L}_{\Lambda}T_{V_{2}}(-)$ 给出。

把 $\mathcal{C}=D_{\mathrm{lis}}(\Bun_{G},\Lambda)$ 看作 $\Rep_{W^{I}_{E}}(\Lambda)$-线性幺半范畴。令 $\End(\mathcal{C})^{BW^{I}_{E}}$ 为带有态射 $W^{I}_{E}\to\Aut(F)$ 的函子 $F\in\End(\mathcal{C})$ 构成的范畴。这仍是幺半范畴。若 $f:*\to[*/(W^{I}_{E})]$，则有态射 $W^{I}_{E}\to\Aut(f)$。这说明从 $\Bun_{G}\to\Bun_{G}\times[*/(W^{I}_{E})]$ 的拉回函子具有这样的自同构；于是从前面的 $T_{V}$ 得到一个元素，仍记作
\[
T_{V}\in\End(\mathcal{C})^{BW^{I}_{E}}.
\]
前三条性质证明了以下结果。

\begin{thm}{5.7}
对应 $V\mapsto T_{V}$ 对 $I$ 具有函子性地定义一个幺半的、$\Rep_{W^{I}_{E}}(\Lambda)$-线性函子
\[
\Rep_{\Lambda}({}^{L}G)^{I}\longrightarrow\End(\mathcal{C})^{BW^{I}_{E}}.
\]
\label{thm:a3-5-7}
\end{thm}

这正是我们用来构造两个中心之间态射的数据。

\subsection{两个中心之间的态射}
\subsubsection{两个中心之间的态射}
设 $\Lambda\in\{\mathbb{Z}_{\ell}[\sqrt{q}],\mathbb{Q}_{\ell}[\sqrt{q}]\}$。对 $\mathbb{Z}_{\ell}[\sqrt{q}]$ 的情形，还假设 $\ell\gg0$ 且满足一个显式的界（「$\widehat{G}$ 的极好素数」；$\mathrm{GL}_{n}$ 时为所有素数，经典群时 $\ell\ne2$，显式定义见 [9, Chapter VII] 的引言）。

\ysrc{156}
把 $\LocSys_{\widehat{G}}$ 从 $\mathbb{Z}_{\ell}$ 基变换到 $\Lambda$。考虑 $\Lambda$-代数
\[
\mathfrak{Z}^{spec}(G,\Lambda)=\mathcal{O}(Z^{1}(W_{E},\widehat{G}))^{\widehat{G}}
\]
（谱稳定 Bernstein 中心），它是 $\LocSys_{\widehat{G}}$ 上凝聚层范畴的中心；以及范畴中心
\[
\mathfrak{Z}^{geo}(G,\Lambda)=\mathfrak{Z}\bigl(D_{\mathrm{lis}}(\Bun_{G},\Lambda)\bigr)
\]
（几何稳定 Bernstein 中心）。本节说明如何由定理 5.7 构造态射
\[
\mathfrak{Z}^{spec}(G,\Lambda)\longrightarrow\mathfrak{Z}^{geo}(G,\Lambda).
\]
由于 $D(G(E),\Lambda)$ 是 $D_{\mathrm{lis}}(\Bun_{G},\Lambda)$ 的直因子，存在态射
\[
\mathfrak{Z}^{geo}(G,\Lambda)\longrightarrow\mathfrak{Z}(G(E),\Lambda)
\]
指向通常的 Bernstein 中心。与前一态射复合，得到态射
\[
\mathfrak{Z}^{spec}(G,\Lambda)\longrightarrow\mathfrak{Z}(G(E),\Lambda).
\]
利用定理 5.3，可以推出前面提到的半单参数构造 $\pi\mapsto\varphi_{\pi}$，见定理 5.10。

\subsubsection{游历代数}
这里在 $\mathbb{Z}_{\ell}$ 上工作。固定 $W_{E}$ 的野惯性的一个开子群 $P$，使它在 $\widehat{G}$ 上平凡作用。考虑开闭子概形
\[
Z^{1}(W_{E}/P,\widehat{G})\subset Z^{1}(W_{E},\widehat{G}).
\]
定理 5.1 的证明表明，可以把 $W_{E}/P$ 换成一个有限型的离散子群 $W$（本质上是把 $I_{E}/P_{E}=\widehat{\mathbb{Z}}{}^{p}(1)$ 换成 $\mathbb{Z}[\tfrac{1}{q}](1)$），使得
\[
Z^{1}(W_{E}/P,\widehat{G})=Z^{1}(W,\widehat{G}).
\]
为使这一点明确，考虑小范畴 $\mathfrak{F}$，其对象为二元组 $(n,F_{n}\to W)$，其中 $n\in\mathbb{N}_{\ge1}$，$F_{n}$ 是 $n$ 元自由群，$F_{n}\to W$ 是态射。$(n,F_{n}\to W)$ 与 $(m,F_{m}\to W)$ 之间的态射，由使图表
\[
\begin{tikzcd}[column sep=large,row sep=small]
F_{n} \arrow[rr] \arrow[dr] & & W \\
& F_{m} \arrow[ur] &
\end{tikzcd}
\]
交换的态射 $F_{n}\to F_{m}$ 给出。于是显然有带 $\widehat{G}$ 的代数作用的 $\mathbb{Z}_{\ell}$-代数同构
\[
\varinjlim_{(n,F_{n})\in\mathfrak{F}}\mathcal{O}(Z^{1}(F_{n},\widehat{G}))\xrightarrow{\ \sim\ }\mathcal{O}(Z^{1}(W,\widehat{G})).
\]
令
\[
\mathrm{Exc}(W,\widehat{G}):=\varinjlim_{(n,F_{n})\in\mathfrak{F}}\mathcal{O}(Z^{1}(F_{n},\widehat{G}))^{\widehat{G}}.
\]
范畴 $\mathfrak{F}$ 并非余滤的，而只是筛余的（以它为指标的余极限与有限积交换，但不与有限极限交换）；而态射
\[
\mathrm{Exc}(W,\widehat{G})\longrightarrow\mathcal{O}(Z^{1}(W,\widehat{G}))^{\widehat{G}}
\]
先验地只是将 $\ell$ 可逆化之后的同构，因为那时取 $\widehat{G}$-不变量是正合函子。Haboush 关于 $\widehat{G}$-不变量的定理表明，这是 $\mathbb{Z}_{\ell}$-代数之间的泛同胚。尽管如此，我们证明下述结果

\ysrc{157}
用到代数群 $\widehat{G}$ 的模表示理论的结果与方法。这是一个更重要结果的推论，我们将在后文解释。

\begin{thm}{5.8}
若 $\ell$ 是极好素数，则 $\mathrm{Exc}(W,\widehat{G})\xrightarrow{\ \sim\ }\mathcal{O}(Z^{1}(W,\widehat{G}))^{\widehat{G}}$。
\label{thm:a3-5-8}
\end{thm}

\subsubsection{游历算子与中心（[9, Chapter VIII.4]）}
固定 $W_{E}$ 的一个有限商 $Q$，使在 $\widehat{G}$ 上的作用经由它分解。这里在任意 $\mathbb{Z}_{\ell}$-代数 $\Lambda$ 上工作。

\begin{thm}{5.9}
设 $\mathcal{C}$ 是任意 $\Lambda$-线性的幂等完备范畴。假设对有限集 $I$ 函子性地给定一个幺半的 $\Rep_{\Lambda}Q^{I}$-线性函子
\[
\Rep_{\Lambda}(\widehat{G}\rtimes Q)^{I}\longrightarrow\End(\mathcal{C})^{BW^{I}}.
\]
则可以构造态射
\[
\mathrm{Exc}(W,\widehat{G})\longrightarrow\mathfrak{Z}(\mathcal{C}),
\]
其中 $\mathfrak{Z}(\mathcal{C})=\End(\mathrm{Id}_{\mathcal{C}})$ 是 $\mathcal{C}$ 的 Bernstein 中心。
\label{thm:a3-5-9}
\end{thm}

事实上，给定 $\mathfrak{F}$ 的一个元素，可以把 $\mathcal{O}(Z^{1}(F_{n},\widehat{G}))^{\widehat{G}}$ 改写成
\[
\mathcal{O}\bigl((\widehat{G})^{n}\sslash\widehat{G}\bigr),
\]
其中 $\widehat{G}$ 在 $(\widehat{G})^{n}$ 上的作用是扭对角作用
\[
g.(g_{1},\dots,g_{n})=(gg_{1}g^{-\tau_{1}},\dots,gg_{n}g^{-\tau_{n}}),
\]
只要 $F_{n}\to W$ 由 $(\tau_{1},\dots,\tau_{n})$ 给出。设 $A$ 是任意 $\Lambda$-代数。给出态射
\[
\mathcal{O}\bigl((\widehat{G})^{n}\sslash\widehat{G}\bigr)\longrightarrow A
\]
等同于给出态射
\[
\mathcal{O}\bigl((\widehat{G}\rtimes Q)^{n}\sslash\widehat{G}\bigr)\longrightarrow\Map(W^{n},A),
\]
它在 $\mathcal{O}(Q^{n})\to\Map(W^{n},A)$ 上线性。现在再添一个变量来改写：沿 $(g_{1},\dots,g_{n})\mapsto(1,g_{1},\dots,g_{n})$ 的拉回是态射 $\mathcal{O}((\widehat{G}\rtimes Q)^{n+1})\to\mathcal{O}((\widehat{G}\rtimes Q)^{n})$，它诱导同构
\[
\mathcal{O}\bigl(\widehat{G}\backslash(\widehat{G}\rtimes Q)^{n+1}/\widehat{G}\bigr)\otimes_{\mathcal{O}(Q^{n+1})}\mathcal{O}(Q^{n})\xrightarrow{\ \sim\ }\mathcal{O}\bigl((\widehat{G}\rtimes Q)^{n}\sslash\widehat{G}\bigr).
\]
于是问题归结为给出态射
\[
\mathcal{O}\bigl(\widehat{G}\backslash(\widehat{G}\rtimes Q)^{n+1}/\widehat{G}\bigr)\longrightarrow\Map(W^{n+1},A),
\]
它在 $\mathcal{O}(Q^{n+1})$ 上线性。

记 $V\mapsto T_{V}$ 为定理 5.9 中的幺半函子。现在假设给定如 [17] 中的五元组 $\mathcal{D}=(I,V,\alpha,\beta,\gamma)$，其中
\begin{enumerate}[(1)]
\item $I$ 是有限集；
\item $V\in\Rep_{\Lambda}(\widehat{G}\rtimes Q)^{I}$；
\item $\gamma\in W^{I}$；
\item $\mathbf{1}\xrightarrow{\alpha}V|_{\widehat{G}}$（经由 $\widehat{G}\subset\widehat{G}^{I}\subset(\widehat{G}\rtimes Q)^{I}$ 的对角限制）；
\item $V|_{\widehat{G}}\xrightarrow{\beta}\mathbf{1}$。
\end{enumerate}
